\mxunit{automatismes-ch14}{Bilan annuel et consolidation des acquis de Seconde}

\begin{fluencysheet}{lun.~14~juin~2027}{Bilan annuel et consolidation}
\mxfq{Échauffement · Flash}{GEO}{$A\pt{-1;4}$ et $B\pt{3;1}$. Quelles sont les coordonnées de $\vect{AB}$\,?}{$\vc{4}{-3}$.}{}
\mxfq{Cœur · Méthode}{GEO}{Donner l'équation réduite de la droite $(AB)$ avec $A\pt{0;1}$ et $B\pt{2;5}$.}{Pente $2$\,; $y=2x+1$.}{}
\mxfq{Cœur + · Méthode}{REP}{$A\pt{1;0}$, $B\pt{4;1}$, $C\pt{5;4}$. Déterminer $D$ pour que $ABCD$ soit un parallélogramme.}{$D=A+C-B=\pt{2;3}$\,; contrôle~: $\vect{AB}=\vect{DC}=\vc31$.}{}
\mxfq{Retour · Flash}{CAL}{Calculer $-5+8\times(-2)$.}{$-5-16=-21$.}{Priorité de la multiplication.}
\mxfq{Fil rouge · Énoncé → formule}{REG}{L'aire d'un carré de côté $x$ est une fonction $A$ de $x$. Donner $A(x)$, un tableau de valeurs pour $x=1\,;\,2\,;\,3$, et dire si $A$ est proportionnelle à $x$.}{$A(x)=x^2$\,; $1\,;\,4\,;\,9$\,; non proportionnelle~: $A(2)\neq2A(1)$.}{}
\end{fluencysheet}

\begin{fluencysheet}{mar.~15~juin~2027}{Bilan annuel et consolidation}
\mxfq{Échauffement · Flash}{NUM}{Écrire $\num{0,00045}$ en notation scientifique.}{$\num{4,5}\times10^{-4}$.}{}
\mxfq{Cœur · Flash}{NUM}{Déterminer $\intfo{-2}{3}\,\cap\,\intof{0}{5}$.}{$\intoo{0}{3}$.}{}
\mxfq{Cœur + · Méthode}{EQ}{Résoudre $|x-1|<4$ et donner $S$ sous forme d'intervalle.}{$-4<x-1<4$~: $S=\intoo{-3}{5}$.}{}
\mxfq{Retour · Développer}{DEV}{Développer et réduire $(2x-1)^2$.}{$4x^2-4x+1$.}{}
\mxfq{Fil rouge · Flash}{GEO}{Un cylindre a pour rayon $3$ cm et pour hauteur $5$ cm. Calculer son volume exact.}{$V=\pi\times3^2\times5=45\pi$ cm$^3$.}{}
\end{fluencysheet}

\begin{fluencysheet}{jeu.~17~juin~2027}{Bilan annuel et consolidation}
\mxfq{Échauffement · Flash}{IMG}{$f(x)=x^2-4x$. Calculer $f(3)$.}{$9-12=-3$.}{}
\mxfq{Cœur · Flash}{EQ}{Résoudre $x^2=9$ puis $|x|=2$.}{$x=\pm3$\,; $x=\pm2$.}{}
\mxfq{Cœur + · Méthode}{REG}{$f(x)=2x-3$ et $g(x)=x+1$. Résoudre $f(x)<g(x)$.}{$2x-3<x+1\Leftrightarrow x<4$~: $S=\intoo{-\infty}{4}$.}{}
\mxfq{Retour · Flash}{QP}{Un robinet débite $18$ L en $4$ min. Quel volume en $10$ min\,?}{$\num{4,5}\times10=45$ L.}{}
\mxfq{Fil rouge · Lecture de code}{ALG}{Quelle est la valeur finale de \code{n}\,?\mxpseudo{n = 0\\ while n*n \textless{} 30:\\ \mxtab n = n + 1}}{$n=6$ car $5^2=25<30$ mais $6^2=36\geqslant30$.}{Plus petit entier naturel dont le carré est supérieur ou égal à $30$.}
\end{fluencysheet}

\begin{fluencysheet}{lun.~21~juin~2027}{Bilan annuel et consolidation}
\mxfq{Échauffement · Flash}{STA}{Quelle est la moyenne de $5,\,7,\,9,\,11$\,?}{$8$.}{}
\mxfq{Cœur · Méthode}{STA}{Déterminer $Q_1$, la médiane et $Q_3$ de la série $1\,;\,3\,;\,4\,;\,6\,;\,8\,;\,9\,;\,11$.}{$n=7$~: $Q_1$ rang $\lceil7/4\rceil=2$~: $3$\,; médiane (rang $4$)~: $6$\,; $Q_3$ rang $\lceil21/4\rceil=6$~: $9$.}{}
\mxfq{Cœur + · Vrai/faux}{STA}{Vrai ou faux\,? Justifier. Ajouter $10$ à chaque valeur d'une série laisse son écart type inchangé.}{Vrai. Les écarts à la moyenne ne changent pas.}{}
\mxfq{Retour · Réflexion}{IMG}{Combien d'antécédents ont $-1$ et $4$ pour $f(x)=x^2$\,?}{$-1$~: aucun\,; $4$~: deux ($-2$ et $2$).}{}
\mxfq{Fil rouge · Prouver}{LIT}{Montrer que pour tout réel $x$~: $(x-1)(x+1)+1=x^2$.}{$(x-1)(x+1)+1=x^2-1+1=x^2$.}{}
\end{fluencysheet}

\begin{fluencysheet}{mar.~22~juin~2027}{Bilan annuel et consolidation}
\mxfq{Échauffement · Flash}{DEV}{Développer et réduire $(2x-3)(x+4)$.}{$2x^2+8x-3x-12=2x^2+5x-12$.}{}
\mxfq{Cœur · Flash}{EQ}{Résoudre $3x-7\leqslant2x+5$.}{$x\leqslant12$\,; $S=\intof{-\infty}{12}$.}{}
\mxfq{Cœur + · Méthode}{EQ}{Résoudre $(x-1)(x+4)=0$ puis $x^2=5$.}{$x=1$ ou $x=-4$\,; $x=\sqrt5$ ou $x=-\sqrt5$.}{}
\mxfq{Retour · Réflexion}{PNP}{Abonnement~: $5$ € + $3$ € par entrée. Le prix est-il proportionnel au nombre d'entrées\,?}{Non~: $1$ entrée $\mapsto8$ €, $2$ entrées $\mapsto11\neq16$ €.}{}
\mxfq{Fil rouge · Flash}{SOL}{Un cône a pour rayon $3$ cm et pour hauteur $4$ cm. Calculer sa génératrice et son volume exact.}{Génératrice $\sqrt{3^2+4^2}=5$ cm\,; $V=\dfrac13\pi\times9\times4=12\pi$ cm$^3$.}{}
\end{fluencysheet}

\begin{fluencysheet}{jeu.~24~juin~2027}{Bilan annuel et consolidation}
\mxfq{Échauffement · Flash}{PRO}{Quel coefficient multiplicateur correspond à une hausse de $\pc{8}$\,?}{$\num{1,08}$.}{}
\mxfq{Cœur · Méthode}{PRO}{Une hausse de $\pc{25}$ est suivie d'une baisse de $\pc{20}$. Quel est le taux global\,?}{$\num{1,25}\times\num{0,8}=1$~: évolution nulle ($\pc{0}$).}{}
\mxfq{Cœur + · Flash}{PRO}{Un prix passe de $\num{5000}$ à $\num{5750}$ Ar. Quel est le taux d'évolution\,?}{$\dfrac{750}{\num{5000}}=\pc{15}$.}{}
\mxfq{Retour · Reconnaître}{PER}{Quelle est la section d'un cylindre par un plan contenant son axe\,? Par un plan perpendiculaire à son axe\,?}{Un rectangle\,; un disque.}{}
\mxfq{Fil rouge · Problème}{EQ}{Le triple d'un nombre, diminué de $4$, est égal au double de ce nombre augmenté de $5$. Écrire une équation puis trouver le nombre.}{$3x-4=2x+5$, d'où $x=9$.}{Vérification~: $3\times9-4=23=2\times9+5$.}
\end{fluencysheet}

\begin{fluencysheet}{lun.~28~juin~2027}{Bilan annuel et consolidation}
\mxfq{Échauffement · Flash}{MOD}{Quel est le sens de variation de $x\mapsto-3x+2$\,?}{Strictement décroissante (pente $-3<0$).}{}
\mxfq{Cœur · Méthode}{REG}{Quels sont le minimum et le maximum de $f(x)=x^2$ sur $\intff{-3}{1}$\,?}{Minimum $0$ (en $0$)\,; maximum $9$ (en $-3$)\,; $f(1)=1<9$.}{}
\mxfq{Cœur + · Méthode}{CMP}{Pour $x>0$, comparer $\dfrac1x$ et $\dfrac1{x+1}$.}{$0<x<x+1$ et l'inverse est décroissante sur $\intoo{0}{+\infty}$~: $\dfrac1x>\dfrac1{x+1}$.}{}
\mxfq{Retour · Programme de construction}{CON}{Décrire en trois étapes la construction du symétrique $A'$ de $A$ par rapport à une droite $(d)$.}{1. Perpendiculaire à $(d)$ passant par $A$, coupant $(d)$ en $H$. 2. Reporter au compas $HA$ de l'autre côté de $(d)$. 3. $A'$ est le point obtenu~: $H$ est le milieu de $[AA']$.}{}
\mxfq{Fil rouge · Méthode}{DEM}{Un point $I$ vérifie $IA=IB$. Où se situe-t-il\,? Quelle propriété permet de l'affirmer\,?}{Sur la médiatrice de $[AB]$ (caractérisation de la médiatrice).}{}
\end{fluencysheet}

\begin{fluencysheet}{mar.~29~juin~2027}{Bilan annuel et consolidation}
\mxfq{Échauffement · Flash}{PRB}{$P(A)=\num{0,28}$. Calculer $P(\overline A)$.}{$\num{0,72}$.}{}
\mxfq{Cœur · Méthode}{PRB}{$100$ personnes dont $60$ femmes\,; $6$ femmes et $4$ hommes sont gauchers. Calculer $P_F(G)$ et $P_G(F)$.}{$P_F(G)=\dfrac6{60}=\num{0,1}$\,; $P_G(F)=\dfrac6{10}=\num{0,6}$.}{}
\mxfq{Cœur + · Méthode}{PRB}{$P(A)=\num{0,4}$, $P_A(B)=\num{0,5}$, $P_{\overline A}(B)=\num{0,2}$. Calculer $P(B)$.}{$\num{0,4}\times\num{0,5}+\num{0,6}\times\num{0,2}=\num{0,2}+\num{0,12}=\num{0,32}$.}{}
\mxfq{Retour · Problème}{MOD}{Volume d'un cube d'arête $x$~: $V(x)=x^3$. Calculer $V(3)$ et trouver $x$ tel que $V(x)=64$.}{$V(3)=27$\,; $x^3=64$, $x=4$.}{}
\mxfq{Fil rouge · Flash}{REP}{Dans quel quadrant (numérotés de $\mathrm{I}$ à $\mathrm{IV}$ dans le sens trigonométrique) se trouve $M\pt{a;b}$ avec $a<0$ et $b>0$\,?}{Quadrant $\mathrm{II}$ (en haut à gauche).}{}
\end{fluencysheet}

\begin{fluencysheet}{jeu.~1\textsuperscript{er}~juil.~2027}{Bilan annuel et consolidation}
\mxfq{Échauffement · Flash}{GEO}{Quelle est la longueur de l'hypoténuse d'un triangle rectangle dont les côtés de l'angle droit mesurent $5$ et $12$\,?}{$\sqrt{25+144}=13$.}{}
\mxfq{Cœur · Flash}{REP}{Quelles sont les coordonnées du milieu de $[AB]$ avec $A\pt{-3;2}$ et $B\pt{5;-4}$\,?}{$\left(\dfrac{-3+5}{2};\dfrac{2-4}{2}\right)=\pt{1;-1}$.}{}
\mxfq{Cœur + · Rédiger}{DEM}{Montrer que $A\pt{0;0}$, $B\pt{2;1}$, $C\pt{-1;2}$ forment un triangle rectangle isocèle en $A$.}{$AB^2=5$, $AC^2=5$\,; $BC^2=9+1=10=AB^2+AC^2$~: réciproque de Pythagore.}{}
\mxfq{Retour · Flash}{NUM}{Donner une valeur approchée de $\dfrac13$ à $10^{-2}$ près. Ce nombre est-il décimal\,?}{$\num{0,33}$. Non~: l'écriture décimale de $\frac13$ est illimitée ($\num{0,333}\ldots$).}{}
\mxfq{Fil rouge · Problème}{DIV}{Quel est le plus petit entier qui donne reste $1$ dans la division par $4$ et par $6$, supérieur à $1$\,?}{$1+\ppcm\pt{4;6}=13$.}{}
\end{fluencysheet}
